problem:  
Let \( M_i^n \subset S^{n+k} \) be a sequence of smooth, compact, **minimal submanifolds** in the unit sphere satisfying  
\[
\|A_{M_i}\|_{L^n(M_i)} \le C, \qquad \mathrm{Vol}(M_i)\le V_0.
\]  
Assume (after subsequence) the varifold convergence \(M_i \rightharpoonup M_\infty\) to a stationary integral varifold \(V_\infty\), with smooth convergence away from finitely many points \( \Sigma = \{p_1, \dots, p_\ell\} \subset S^{n+k}\).

Consider the **Jacobi operator**
\[
L_{M_i} = \Delta_{M_i}^{\perp} + |A_{M_i}|^2 + n,
\]
acting on smooth normal sections of \(M_i\), and denote by  
\[
\operatorname{ind}(M_i) = \#\{\text{negative eigenvalues of }L_{M_i}\}
\]
the **Morse index** of \(M_i\).

**Question (Morse index compactness under critical curvature bound):**  
Under the above \(L^n\)-curvature and volume bounds, is the Morse index \(\operatorname{ind}(M_i)\) **upper semicontinuous** under varifold convergence, up to finitely many bubble contributions?  

That is, can one find finitely many complete minimal submanifolds (bubbles) \(\Sigma_j^\ast \subset \mathbb{R}^{n+k}\), arising as blow-up limits at each \(p_j \in \Sigma\), such that
\[
\limsup_{i\to\infty}\operatorname{ind}(M_i) \le \operatorname{ind}(M_\infty) + \sum_{j=1}^\ell \operatorname{ind}(\Sigma_j^\ast)?
\]

Equivalently, does **index quantization** hold at the critical curvature threshold—so that all loss of index along a convergent sequence of \(L^n\)-bounded minimal submanifolds is captured exactly by bubbling phenomena at singular points?

If not, can one construct counterexamples in which the Morse index diverges even though curvature and volume stay uniformly bounded?

Why is it a "good" problem:  
This problem directly targets the unresolved interface between **analytic stability theory** and **geometric compactness** for minimal submanifolds. The Morse index measures how many independent destabilizing directions are available under area variations; understanding its behavior under geometric degeneration is fundamental to the structure of the moduli space of minimal submanifolds.

The question tests whether analytic stability behaves continuously under **critical curvature compactness**. It unifies three deep themes:

1. **Geometric measure theory:** convergence and bubbling for minimal submanifolds with curvature concentration.
2. **Spectral theory of elliptic operators:** semicontinuity of the number of negative eigenvalues in degenerating families.
3. **Analytic thresholds:** the \(L^n\) bound is precisely the minimal integrability ensuring weak compactness of curvature and scale invariance of curvature energy.

If index quantization holds, it would extend the bubble-tree structure for minimal submanifolds by incorporating not only geometric but also spectral data, yielding a bridge between geometric compactness and spectral stability—a result paralleling index quantization in harmonic map and Yang–Mills theories.

The statement is simple—how does the Morse index behave in \(L^n\)-bounded minimal geometry?—yet profound, striking at the analytic core of compactness and stability in differential geometry.